\section{Removing an unbounded prime-to-\texorpdfstring{$p$}{p} kernel}\label{sec:kernel}

We now turn the finite based lists on $p$-subgroups into a finite
Morita list for block summands of the whole crossed order.  Its
grading group may have unbounded order.  The bounded quantity is
its index over its largest normal $p'$-subgroup.

The proof refines the field argument of
\cite[Lemma~8.7]{OpenAIFieldDonovan2026} through the
Clifford-theoretic matrix corners of
\cite{KulshammerClifford1990,Kulshammer1995}.  Over $\cO$, the
principal-unit argument uses unique prime-to-$p$ divisibility.
Removing a matrix factor creates a second scalar error; determinant
normalization places it in $\T$, so the original based Sylow
restriction is recovered exactly.  After these two steps, the
remaining crossed systems have bounded grading groups and finite
fibres under restriction.

\subsection{Central units and the restriction kernel}

The principal-unit argument follows
\cite[Lemma~4.4]{EiseleReduction2021}.  Sylow restriction will control
the part for which the grading group may have $p$-torsion.

\begin{lemma}[Units of the centre]\label{lem:center-units}
Let $R$ be an indecomposable $\cO$-order and put $Z_R=Z(R)$.
Then $Z_R$ is local with residue field $k$, and
\begin{equation}\label{eq:center-units}
 Z_R^\times=\T\times U_R,
 \qquad U_R=1+\rad Z_R.
\end{equation}
This decomposition is invariant under every $\cO$-algebra
automorphism of $R$.  If $n$ is prime to $p$, the map
$u\mapsto u^n$ is an automorphism of $U_R$.
\end{lemma}

\begin{proof}
The finite commutative $\cO$-algebra $Z_R$ is a product of
complete local algebras, by henselianity of $\cO$.  More than
one factor would give a nontrivial central idempotent of $R$.
It is therefore local.  Its residue field is a finite extension
of the algebraically closed field $k$, hence is $k$.
The residue map on units is split by the scalar Teichm\"uller
units, and its kernel is $U_R$, proving
Equation~\eqref{eq:center-units}.  Both factors are canonical
and are preserved by the stated automorphisms; these act
trivially on $\T$.

For $u\in U_R$, the polynomial $X^n-u$ has the root $1$ modulo
the maximal ideal, with derivative $n$ a unit.  Hensel's lemma
gives a unique root in $U_R$.  This proves the last assertion.
\end{proof}

\begin{samepage}
\begin{lemma}[Finite kernel of Sylow restriction]
\label{lem:restriction}
Let a finite group $Q$ act by $\cO$-algebra automorphisms on
$Z_R$, and let $S$ be a Sylow $p$-subgroup of $Q$.  Then
\[
 \res:H^2(Q,Z_R^\times)\longrightarrow H^2(S,Z_R^\times)
\]
has finite kernel.
\end{lemma}
\end{samepage}

\begin{proof}
The two factors in Equation~\eqref{eq:center-units} play different
roles: restriction is injective on the principal-unit part, while
the entire Teichm\"uller cohomology group is finite.  On $U_R$ the composite
$\cor\circ\res$ is multiplication by $[Q:S]$.  This is an
automorphism on $U_R$, by Lemma~\ref{lem:center-units}, and
therefore on its cohomology.  Restriction is consequently
injective on $H^2(Q,U_R)$.

On $\T$, the action is trivial and $H^2(Q,\T)$ is finite.
Indeed, let $d=|Q|$.  The $d$th-power map on $\T$ is onto
and has the finite kernel $\mu_d(\cO)\cap\T$.  The exact
sequence in cohomology and the annihilation of positive
cohomology by $d$ show that $H^2(Q,\T)$ is a quotient of
\[
 H^2(Q,\mu_d(\cO)\cap\T),
\]
which is finite because both the group and the coefficient
module are finite.  Combining the two factors proves the lemma.
\end{proof}

\subsection{The integral extension theorem}

\begin{theorem}[Integral kernel removal]\label{thm:kernel}
Fix an indecomposable basic $\cO$-order $R$ Morita equivalent
to a finite-group block, and a finite subgroup
$\mathcal F\leq\Out_{\cO}(R)$.  Consider crossed orders on
$R$ by finite groups $Q$ with the following properties:
\begin{enumerate}
\item the indices $[Q:O_{p'}(Q)]$ are bounded;
\item every outer action has image in $\mathcal F$;
\item for each possible abstract $p$-group, the restrictions
to subgroups of that type have finitely many based graded
isomorphism classes.
\end{enumerate}
Then the block summands of these crossed orders have only
finitely many $\cO$-linear Morita equivalence classes.
\end{theorem}

\begin{proof}
Write $A=\bigoplus_{q\in Q}Ru_q$ for one crossed order and set
\begin{equation}\label{eq:large-kernel}
 K=O_{p'}(Q)\cap\ker(Q\longrightarrow\mathcal F).
\end{equation}
It is a normal $p'$-subgroup, and
$[Q:K]\leq[Q:O_{p'}(Q)]|\mathcal F|$ is bounded.
We isolate its twisted group algebra, pass to matrix corners,
verify their Sylow restrictions, and finally count central factor
systems on the bounded quotient.

\emph{Scalar factors on $K$.}
Since $K$ acts trivially modulo inner automorphisms, change the
$K$-degree units so that they centralize $R$.  Their factors
belong to $Z_R^\times$ and form an ordinary $2$-cocycle for
the trivial $K$-action.  Multiplication by $|K|$ is invertible
on $U_R$, so $H^i(K,U_R)=0$ for $i>0$.  Removing the
$U_R$-component by a central change of units gives units $v_k$
with
\[
 v_kv_l=\beta(k,l)v_{kl},\qquad \beta(k,l)\in\T.
\]
Their $\cO$-span is a twisted group algebra
$E=\cO_\beta K$, and the restriction to $K$ is
$R\otimes_{\cO}E$.

Every homogeneous unit of $A$ normalizes $E$.  To verify this,
write
\[
 u_qv_ku_q^{-1}=z_kv_{qkq^{-1}},\qquad z_k\in Z_R^\times.
\]
The coefficient is central because both sides centralize $R$.
Comparing the products for $k,l\in K$ shows that their
$U_R$-components define a homomorphism $K\to U_R$: the
original factors, their conjugates, and the factors in the
new $K$-degrees all lie in $\T$.  Such a homomorphism is
trivial, since multiplication by $|K|$ is invertible on
$U_R$.  Hence every $z_k$ belongs to $\T$, proving
normalization of $E$.

\emph{The matrix factors of $E$.}
The finitely many values of $\beta$ generate a finite
$p'$-subgroup of $\T$.  The usual central extension defined
by $\beta$ is therefore a finite $p'$-group, and $E$ is its
character summand over $\cO$.  A finite $p'$-group algebra
over $\cO$ is a product of full matrix algebras: its reduction
is split semisimple, and the matrix units lift over the
complete ring $\cO$.  Each lifted primitive corner has rank
one and is $\cO$.  We obtain
\begin{equation}\label{eq:matrix-E}
 E\cong\prod_j M_{n_j}(\cO).
\end{equation}
Every $n_j$ is prime to $p$.  Indeed these are ordinary
irreducible character degrees in the indicated character
sector of a finite $p'$-group; each such degree divides the
group order.

Let $e_j$ denote the identity of one matrix factor.  The
$Q$-orbits on these idempotents give central summands of
$A$.  In each orbit summand a single $e_j$ is full, because
its homogeneous conjugates sum to that orbit idempotent.
Let $Q_j^{\mathrm{pre}}$ be its stabilizer.  It contains
$K$, and the $e_j$-corner, with its grading coarsened by
$K$, is crossed over
\[
 R\otimes_{\cO}M_n(\cO),\qquad n=n_j,
\]
by the group $Q_j=Q_j^{\mathrm{pre}}/K$.  The order of $Q_j$
is bounded by $[Q:K]$.

In each $Q_j$-degree choose a unit $e_ju_q$ inherited from
an original $Q$-degree.  Its conjugation preserves both $R$
and $M_n(\cO)$, by the normalization of $E$ proved above.
Every $\cO$-algebra automorphism of
$M_n(\cO)$ is inner.  For example, its images of the standard
matrix idempotents split $\cO^n$ into free rank-one summands;
choosing compatible basis vectors for the matrix units
constructs an implementing matrix.  Correct the homogeneous
units by these matrices so that they centralize the matrix
factor.  Their multiplication factors then lie in its
centralizer inside $R\otimes_{\cO}M_n(\cO)$, namely $R$.
The orbit corner is consequently
\[
 M_n(\cO)\otimes_{\cO}C_j,
\]
where $C_j$ is a $Q_j$-crossed order on $R$ with the same
outer action on $R$.  Equivalently, $C_j$ is the centralizer
of the matrix factor in the orbit corner.

This passage replaces the unbounded group $Q$ by the bounded group
$Q_j$.  It has not yet supplied a finite list: the matrix corrections
may change the multiplication factors of a restriction.  We next
recover those factors.  The matrix size need not be bounded; its
prime-to-$p$ property is what the recovery uses.

\emph{Preservation of the based Sylow restriction.}
Let $S$ be a Sylow $p$-subgroup of $Q_j$.  Schur--Zassenhaus
in its inverse image gives a $p$-subgroup
$\widetilde S\leq Q_j^{\mathrm{pre}}$ mapping isomorphically
onto $S$.  Use the original units $u_s$ in these degrees,
with original factors $a_{s,t}\in R^\times$.  Their
conjugations on $E_j=M_n(\cO)$ form an honest group action,
since $a_{s,t}$ centralizes $E_j$.

Choose implementing matrices $c_s\in\mathrm{GL}_n(\cO)$,
with $c_1=1$.  We may take $\det c_s=1$.  In fact, since
$p\nmid n$, every unit of $\cO$ has an $n$th root: take a
Teichm\"uller root of its residue and then use Hensel's
lemma on the principal-unit part.  Multiplication by a
scalar thus normalizes the determinant without changing
the implemented automorphism.

Because the automorphisms form a group action, there is a
scalar $\gamma(s,t)\in\cO^\times$ such that
\[
 c_sc_tc_{st}^{-1}=\gamma(s,t)I_n.
\]
Taking determinants gives $\gamma(s,t)^n=1$.  Thus
$\gamma(s,t)\in\T$, since $n$ is prime to $p$.
Associativity shows that $\gamma$ is a normalized
$\T$-valued $2$-cocycle on $S$.

For the corrected units $w_s=c_s^{-1}e_ju_s$, their
actions on $R$ are unchanged.  The formula for their
products is
\begin{equation}\label{eq:corrected-matrix-factors}
 w_sw_t=\gamma(s,t)^{-1}a_{s,t}w_{st}.
\end{equation}
For example, this follows by using
$u_sc_tu_s^{-1}=c_sc_tc_s^{-1}$ in the product on the
left.  Lemma~\ref{lem:teich-cohomology} removes $\gamma$
by scalar rescaling.  Therefore the restriction of $C_j$
to $S$ is based-isomorphic to the original restriction to
$\widetilde S$.  The corrected units generate the corresponding
homogeneous $R$-components of the centralizer $C_j$.  For a fixed
original unit, any two implementing matrices differ by a scalar
unit; the subsequent scalar rescaling also changes only the
generator of its component.  Thus these choices change the
presentation, not the based graded order.  Hypothesis (iii) supplies a finite list
for the restrictions now obtained.

\emph{Crossed systems on a bounded quotient.}
There are finitely many possible abstract groups $Q_j$
and outer homomorphisms $Q_j\to\mathcal F$.  Fix one of
each and choose representatives
$\alpha_q\in\Aut_{\cO}(R)$ of its outer classes.
Changing homogeneous units makes every crossed system
with this outer action use these same automorphisms.
If any factors satisfying the crossed identities exist,
their classes under central changes of units form a
torsor under
\begin{equation}\label{eq:factor-torsor}
 H^2(Q_j,Z_R^\times).
\end{equation}
Indeed, the ratio of two factor systems with these fixed
actions is central by Equation~\eqref{eq:action-law}.
Equation~\eqref{eq:factor-law} says that the ratio is a
$2$-cocycle.  A homogeneous-unit change preserving every
$\alpha_q$ must itself be central, and changes the ratio
by a coboundary.  The action on the centre is a genuine
group action because inner automorphisms act trivially
there.

On a Sylow $p$-subgroup, the based finiteness just proved
gives finitely many restricted torsor classes.  To see
that the fixed representatives cause no extra ambiguity,
observe that a based isomorphism between systems with
the same automorphisms must change their units centrally.
Lemma~\ref{lem:restriction} gives finite fibres for the
restriction map from Equation~\eqref{eq:factor-torsor}: a nonempty
fibre is a coset of its finite kernel.
Thus there are finitely many based crossed orders $C_j$.

Each $C_j$ has only finitely many block summands.  The
orbit corners above are matrix algebras over them, and
the selected $e_j$ was full in the corresponding orbit
summand of $A$.  Hence every block summand of $A$ is
Morita equivalent to one of this finite list of block
summands of the $C_j$.  Neither the number of matrix
factors in Equation~\eqref{eq:matrix-E} nor their sizes
had to be bounded.  This proves the theorem.
\end{proof}

The two scalar normalizations in this proof have different reasons.
On $K$, its prime-to-$p$ order makes principal-unit cohomology vanish.
On the lifted Sylow subgroup, the determinant first forces the error
into $\T$, and only its $\T$-cohomology is used.  Neither step asserts
vanishing of principal-unit cohomology on a $p$-group.
